\section{Conclusion}
\label{sec:conclusion}

We built OPENER for legal NER, a frugal open-world pipeline assembled from ready-to-use pretrained parts and deliberately kept free of any legal-specific training, and evaluated both of its typing heads on four legal corpora in three languages against the standard open-world baselines.

The two heads diverge sharply on legal text. OPENER-ZS remains a usable label-free option on English legal text, but it does not lead the zero-shot field: here it trails GLiNER-M/L and OWNER on average, on gold mentions as well as end-to-end, and its gold AMI collapses cross-lingually as its English-centric prototypes stop matching target-language mentions. OPENER-Sup is almost unaffected, holding up even on German-LER's fine $19$-type schema, and it posts the best end-to-end AMI on average, a statistically significant lead. The reason is geometric: a balanced probe fits hyperplanes precise enough to pull apart legal types whose vectors sit close together, which a single label-name match cannot do without supervision. \rev{The confusion counts make this concrete: the label-free head maps German \textsc{norm} onto \textsc{regulation} $213$ times, the probe $3$ times.} The supervised head is therefore the safer operating point wherever labels exist, and OPENER-ZS is best reserved for languages close to the embedder's own.

End-to-end quality remains bounded first by mention \emph{detection}, whichever typing head runs behind it. \rev{A typing oracle over the detected spans stops at ${\approx}42$ AMI, and most of what detection misses are boundary disagreements rather than mentions it never sees.} Adapting the detector in-language, rather than the embedder, is the natural next step and preserves the frugal decoupled design. Nothing else in the design is tied to legal text beyond the type names it is asked about: legal NER was the demanding case we built it for, and it holds.

\section*{Limitations}
\rev{\paragraph{Scale of the LLM comparison} Our LLM baseline is a $4$-bit Qwen2.5-1.5B, and it is the largest instruction-tuned model our hardware runs: every system in this paper is measured on one consumer GPU with $6$\,GB of VRAM, which a $7$B model saturates in $4$-bit and a larger one does not fit at all. Frontier models are reachable only through paid APIs, and this study has no funding line for inference credits. That constraint is not only budgetary. A hosted model would break the very comparison this paper is built on, since two of our three axes cannot be measured through an API: latency would fold in network round-trips, provider-side batching and queueing, none of which is a property of the model, and energy could not be measured at all, the serving hardware being neither known nor instrumentable by CodeCarbon. Reporting an API model beside locally measured systems would therefore mean either dropping the cost axes or filling them with numbers that are not comparable. We take the honest consequence: we do not know how close a frontier LLM would come to the topline on these corpora, and the oracle-typing ceiling of \autoref{sec:exp:results} is our model-independent answer to how much room is left above the systems we could measure.}

Test splits are capped at their first $1000$ entity-bearing sentences, which affects German-LER only, and sentences carrying no gold entity are excluded, so precision on entity-free text is never measured. OPENER scores are means over three embedder-retraining seeds: the supervised and end-to-end scores are stable (std ${\leq}1.4$ AMI), but the label-free prototypes vary by up to ${\pm}6$ AMI across retrainings, so single-seed zero-shot numbers should be read with care. Inference is deterministic for a fixed embedder, and test-set bootstraps on the released model confirm both the typing-on-gold scores (std ${\pm}0.8$--$1.7$ AMI) and the significance of OPENER-Sup's end-to-end lead, stable across resampling seeds $42/7/123$. More fundamentally, both the embedder and the detector are English-centric and used without any legal- or language-domain adaptation, which is the design under test here and also its main exposure. The detectors' low cross-lingual recall (on German, $18.7\%$ for GLiNER-L and $10.2\%$ for GLiNER-M) caps end-to-end quality and is the first thing a follow-up should address.

\section*{Ethics Statement}
This work runs and evaluates entity recognition on four publicly released legal NER benchmarks (E-NER over SEC/EDGAR filings, an Indian court-judgement set, LeNER-Br, and German-LER), used only for evaluation and under their original licenses. Legal documents contain personal data such as the names of parties, judges and counsel. We introduce no new annotations and release no data, we only run inference over the existing public test splits, and we report no individual-level content, so the study adds no privacy exposure and makes no attempt at re-identification.

Open-world NER is error-prone, and our own results show recall dropping sharply outside English (as low as $10.2\%$ on German). The pipeline is therefore meant as an assistive tool that surfaces candidate entities for human review, not as a component of automated legal decision-making, where a missed or mistyped entity can carry material consequences. A missed entity is silent, so a human in the loop remains necessary, especially cross-lingually. Finally, a central motivation of this study is energy frugality. We report per-sentence latency and energy alongside accuracy precisely to make the environmental cost of deployment explicit, and to argue for models that are cheap enough to run on modest hardware.

The authors used Claude Code (Anthropic) to scaffold and debug the evaluation scripts and to draft and copy-edit this manuscript. The research question, the experimental design, the choice of baselines and datasets, the execution of all experiments and the interpretation of the results are the authors' own. The authors reviewed and edited all content and take full responsibility for it.
